% Modern editorial, markup and figure/layout contributions: CC-BY-SA-4.0.
% Historical text and illustrations: Leonhard Euler (1744), public domain.
% Component scope and upstream attribution: see RIGHTS.md.
\sourcepage{5}
quæsitæ sit transeundum; totidem vero puncta, quot insunt in
æquatione inventa quantitates arbitrariæ, ipsam æquationem determinatam
reddent. Loco punctorum autem, ad curvam quæsitam perfecte determinandam,
adhiberi etiam possunt totidem tangentes quæ curvam quæsitam tangant, \&,
si contactus debeat fieri in dato tangentis puncto, hæc conditio duobus
punctis æquivalebit. Quin etiam in locum punctorum, aliæ quæcunque
conditiones substitui possunt; dummodo eæ ita sint comparatæ, ut per eas
quantitates arbitrariæ in æquatione inventa contentæ determinentur.
Neque vero ante opus est solutionem ad finem perducere, quam ista
dijudicatio suscipiatur; sed infra tradentur certa criteria, quorum ope,
statim, ex illa quantitate variabili quæ maximum minimumve esse debet,
dignosci poterit quæ novæ constantes in æquationem pro curva ingrediantur,
quæ in quæstione non continebantur. Oriuntur autem istæ constantes
arbitrariæ ex gradu differentialium, ad quem æquatio pro curva quæsita
exsurgit; quoti enim gradus prodit æquatio differentialis pro curva quæsita,
tot quantitates arbitrariæ in illa censendæ sunt potestate inesse;
hincque totidem conditionibus opus erit ad curvam determinandam. Idem
vero etiam usu venit in solutione omnium Problematum, quando æquatio
differentialis vel primi vel altioris gradus invenitur; ita ut hinc in
præsenti instituto nulla peculiaris difficultas inesse censenda sit.

\subsection*{Definitio III.}
\originalsection{10}\textit{Methodus maximorum ac minimorum relativa} docet, non
inter omnes omnino curvas eidem abscissæ respondentes, sed inter eas tantum
quæ præscriptam quandam proprietatem communem habeant, eam determinare quæ
maximi minimive proprietate gaudeat.

\subsection*{Corollarium I.}
\originalsection{11}Ad hujusmodi igitur Problemata solvenda, primum, ex
omnibus omnino curvis eidem abscissæ respondentibus eæ sunt
\sourcepage{6}segregandæ, in
quas eadem præscripta proprietas competat; atque tum demum ex his segregatis
ea quæ quæritur debebit definiri.

\subsection*{Corollarium II.}
\originalsection{12}Quanquam autem tali conditione numerus curvarum omnium
ad eandem abscissam relatarum vehementer restringitur, tamen is etiamnum
manebit infinitus. Quin etiam, si non una, sed plures proprietates
præscribantur, quibus omnes curvæ ex quibus quæsita est determinanda debeant
esse præditæ, tamen usque numerus curvarum manebit infinitus.

\subsection*{Corollarium III.}
\originalsection{13}Quo plures itaque proponuntur proprietates, quæ iis
curvis ex quibus quæsitam definiri oportet communes esse debeant, eo magis
numerus curvarum inter quas electio quæsita est instituenda restringetur,
etiamsi maneat infinitus.

\subsection*{Scholion I.}
\originalsection{14}Ex hoc genere, in quo Methodum maximorum \& minimorum
relativam constituimus, initio hujus Seculi, primum a Jacobo Bernoullio in
medium prolatum est famosum illud Problema Isoperimetricum; in quo
quærebatur curva maximi minimive proprietate prædita, non inter omnes
curvas ad eandem abscissam relatas, sed inter eas tantum quæ ejusdem
essent longitudinis; ex quo istæ curvæ, ex quibus quæsitam erui oportebat,
isoperimetræ sunt appellatæ. Ita si, inter omnes curvas eidem abscissæ
respondentes \& longitudine æquales, quæratur ea quæ cum abscissa \& applicata
maximum spatium includat; reperitur quæsito linea circularis satisfacere,
quod quidem jam diu ante inventam hanc methodum Geometris innotuerat, ac
demonstratum erat. At, hoc casu iterum, ex ipsa Problematum natura, novæ
conditiones accedunt; uti in iis, quæ ad Methodum maximorum ac minimorum
absolutam pertinent; \sourcepage{7} quæ ex constantibus arbitrariis, quas solutio inducit,
sunt æstimandæ. Ita in solutione Problematis, quo curva quæritur quæ inter
omnes ejusdem longitudinis maximam comprehendat aream cum abscissa, duæ
constantes novæ ingrediuntur; ex quo, ad Problema determinatum efficiendum,
id ita est proponendum, ut inter omnes curvas ejusdem longitudinis, quæ non
solum eidem abscissæ respondeant, sed etiam per data duo puncta transeant,
quæratur ea quæ ad datam abscissam maximam aream referat. Atque simili
modo evenire potest, ut quatuor puncta, \& plura etiam interdum, pro arbitrio
assumi debeant, quo Problema fiat determinatum; cujus rei dijudicatio ex
ipsa Problematum natura est petenda. Quemadmodum autem, in Problemate
isoperimetrico, omnes curvæ ex quibus quæsitam determinari oportet ejusdem
longitudinis ponuntur; ita loco hujus proprietatis alia quæcunque proponi
potest, quæ omnibus communis esse debeat. Sic jam quæsitæ sunt curvæ maximi
minimive proprietate præditæ, inter omnes eas curvas ad eandem abscissam
relatas tantum, quæ circa eam abscissam conversæ omnes æquales superficies
generent; atque simili modo aliæ quæcunque proprietates proponi possunt.
Deinde etiam, non una, sed plures hujusmodi proprietates præscribi possunt,
quæ omnibus curvis inter quas ea quæ maximum minimumve aliquod contineat
definienda sit communes esse debeant. Ita si quæreretur curva maximi vel
minimi proprietate quapiam prædita, inter omnes curvas eidem abscissæ
respondentes, quæ tam essent omnes inter se longitudine æquales, quam etiam
areas æquales concluderent.

\subsection*{Scholion II.}
\originalsection{15}Propter hoc discrimen inter Methodum maximorum \&
minimorum absolutam ac relativam, tractatio nostra erit bipartita. Primum
scilicet methodum trademus, inter omnes omnino curvas eidem abscissæ
respondentes, eam determinandi quæ maximi minimive proprietate sit prædita.
Deinde vero progrediemur ad ejusmodi Problemata, in quibus curva maximi
minimive \sourcepage{8} proprietate gaudens postulatur, inter omnes curvas quæ unam
pluresve propositas proprietates communes habeant; atque ex numero harum
proprietatum istius tractationis denuo subdivisio orietur. Interim tamen
non opus erit in hac subdivisione longius progredi; cum mox reperiatur
methodus, quotcunque etiam propositæ fuerint proprietates, Problemata facile
resolvendi. Solutiones enim Problematum prima fronte maxime intricatorum
præter opinionem fient perquam expeditæ, ac levi calculo absolvendæ.
